\section*{Chapter \ref{ch:g7:how-a-processor-runs-a-program} --- How a Processor Runs a Program}

\begin{solution}{exo:g7:how-a-processor-runs-a-program:1}
\code{R0} holds 5, then \code{R1} holds 2, then \code{R0} becomes $5 - 2 = 3$:
\code{R0} is 3 and \code{R1} is 2.
\end{solution}

\begin{solution}{exo:g7:how-a-processor-runs-a-program:2}
Each fetch adds 2: after the steps \code{PC} holds 2, 4, 6 and 8 (the fourth
step is the first \code{ADD}, which does not jump).
\end{solution}

\begin{solution}{exo:g7:how-a-processor-runs-a-program:3}
\omoutput{code/grade-7/04-how-a-processor-runs-a-program/outputs/countdown.txt}{The
countdown.}\label{out:g7:how-a-processor-runs-a-program:countdown}

\code{OUT} comes before \code{SUB}, so 3 is printed first; when \code{R0}
reaches 0 the jump is not taken, and 0 is never printed.
\end{solution}

\begin{solution}{exo:g7:how-a-processor-runs-a-program:4}
\omcode{code/grade-7/04-how-a-processor-runs-a-program/exercises/sum_plus_one.toy}{2}{8}{The
sum plus 1.}\label{lst:g7:how-a-processor-runs-a-program:sol-sum-plus-one}

With 3 and 4 it prints 8. There is no \omterm{def:g1:instructions-in-order:instruction}{instruction} to add a number directly in
the table, so the 1 goes first into a \omterm{def:g7:how-a-processor-runs-a-program:register}{register}.
\end{solution}

\begin{solution}{exo:g7:how-a-processor-runs-a-program:5}
\code{JNZ again} jumps back to \code{LOADI R0, 3}, which puts 3 back in
\code{R0} at every turn: \code{R0} never reaches 0, and the \omterm{def:g4:programs-and-operating-system:program}{program} prints 3
forever (2\,000 times before the machine stops it). The label belongs on the
\code{OUT} line, after the two \code{LOADI}, as in
\cref{lst:g7:how-a-processor-runs-a-program:countdown}.
\end{solution}

\begin{solution}{exo:g7:how-a-processor-runs-a-program:6}
After \code{OUT R0} the \omterm{def:g3:inside-the-computer:processor}{processor} fetches the next two cells, 8 and 9. Every
cell the \omterm{def:g4:programs-and-operating-system:program}{program} did not fill holds 0, and the two \omterm{def:g4:memory-and-storage:byte}{bytes} 0, 0 are the
\omterm{def:g1:instructions-in-order:instruction}{instruction} \code{HALT}: the machine stops by luck. If anything were stored
after the \omterm{def:g4:programs-and-operating-system:program}{program}, data or another part of the \omterm{def:g4:programs-and-operating-system:program}{program}, the \omterm{def:g3:inside-the-computer:processor}{processor} would
run it as \omterm{def:g1:instructions-in-order:instruction}{instructions}.
\end{solution}

\begin{solution}{exo:g7:how-a-processor-runs-a-program:7}
\omcode{code/grade-7/04-how-a-processor-runs-a-program/exercises/five_to_one.toy}{2}{7}{From
5 down to 1.}\label{lst:g7:how-a-processor-runs-a-program:sol-five}
\end{solution}

\begin{solution}{exo:g7:how-a-processor-runs-a-program:8}
\code{LOADI} has code 1 and the \omterm{def:g7:how-a-processor-runs-a-program:register}{register} is 1: the first \omterm{def:g4:memory-and-storage:byte}{byte} is
$4 \times 1 + 1 = 5$, \code{00000101}; the second is the number 7,
\code{00000111}.
\end{solution}

\begin{solution}{exo:g7:how-a-processor-runs-a-program:9}
Three steps before the \omterm{def:g4:algorithms-repeat-and-decide:loop}{loop} (\code{IN} and two \code{LOADI}), three per turn
for 4 turns, then \code{OUT} and \code{HALT}: $3 + 12 + 2 = 17$, as the trace
says.
\end{solution}

\begin{solution}{exo:g7:how-a-processor-runs-a-program:10}
$256 / 2 = 128$ \omterm{def:g1:instructions-in-order:instruction}{instructions}. A real \omterm{def:g4:programs-and-operating-system:program}{program} also needs cells for its data,
and the last \omterm{def:g1:instructions-in-order:instruction}{instruction} must be \code{HALT}, so fewer \omterm{def:g1:instructions-in-order:instruction}{instructions} do useful
work.
\end{solution}

\begin{solution}{exo:g7:how-a-processor-runs-a-program:11}
Three steps before the \omterm{def:g4:algorithms-repeat-and-decide:loop}{loop}, three per turn for $n$ turns, two at the end:
$3n + 5$ steps.

\omcode{code/grade-7/04-how-a-processor-runs-a-program/listings/sum_steps.py}{1}{7}{Steps
of \cref{lst:g7:how-a-processor-runs-a-program:sum}.}\label{lst:g7:how-a-processor-runs-a-program:sum-steps}

\omoutput{code/grade-7/04-how-a-processor-runs-a-program/outputs/sum_steps.txt}{From
8 steps for $n = 1$, three more for each extra
turn.}\label{out:g7:how-a-processor-runs-a-program:sum-steps}
\end{solution}

\begin{solution}{exo:g7:how-a-processor-runs-a-program:12}
\code{CMP R0, R1} computes $R0 - R1$; it is negative exactly when \code{R0}
is the smaller one, and then the \omterm{def:g4:programs-and-operating-system:program}{program} jumps to print \code{R1}.

\omcode{code/grade-7/04-how-a-processor-runs-a-program/exercises/larger.toy}{2}{9}{The
larger of two numbers.}\label{lst:g7:how-a-processor-runs-a-program:sol-larger}

With 9 and 14 it prints 14. (Numbers of 128 or more would need the flag
\code{C}, which the next book explains.)
\end{solution}

\begin{solution}{pb:g7:how-a-processor-runs-a-program:1}
\textbf{1.} $6 \times 4 = 6 + 6 + 6 + 6$ (three additions to a first 6, or
four to a total starting at 0); $4 \times 6 = 4 + 4 + 4 + 4 + 4 + 4$, with
more additions.

\textbf{2.} The first four rows read $a$, $b$ and set up \code{R2} = 0 and
\code{R3} = 1 (\code{PC} 0, 2, 4, 6). Then each turn runs \code{ADD},
\code{SUB}, \code{JNZ} at addresses 8, 10 and 12: \code{R2} goes 6, 12, 18,
24 while \code{R1} goes 3, 2, 1, 0. Finally \code{OUT} at 14 and
\code{HALT} at 16.

\textbf{3.} $4 + 3 \times 4 + 2 = 18$ steps; the answer is in \code{R2}, and
the \omterm{def:g4:programs-and-operating-system:program}{program} prints 24.

\textbf{4.} \code{R1} counts the additions that are still to be done.

\textbf{5.} It agrees: 18 steps for $b = 4$.

\omoutput{code/grade-7/04-how-a-processor-runs-a-program/outputs/multiply_steps.txt}{The
steps for $6 \times b$.}\label{out:g7:how-a-processor-runs-a-program:multiply-steps}

\textbf{6.} By 3 each time: one more turn of the \omterm{def:g4:algorithms-repeat-and-decide:loop}{loop}, which is three
\omterm{def:g1:instructions-in-order:instruction}{instructions}.

\textbf{7.} The four \omterm{def:g1:instructions-in-order:instruction}{instructions} before the \omterm{def:g4:algorithms-repeat-and-decide:loop}{loop} and the two after it: 6
steps.

\textbf{8.} $3b + 6$ steps.

\textbf{9.} $3 \times 50 + 6 = 156$ steps for $3 \times 50$, and
$3 \times 3 + 6 = 15$ for $50 \times 3$:

\omoutput{code/grade-7/04-how-a-processor-runs-a-program/outputs/orders.txt}{The
order matters for the time, not for the answer (except on the last
line).}\label{out:g7:how-a-processor-runs-a-program:orders}

\textbf{10.} The smaller one, since the number of turns is $b$.

\textbf{11.} $7 \times 200 = 1\,400$ does not fit in one \omterm{def:g4:memory-and-storage:byte}{byte}. The machine
prints 120 instead: it kept only what fits in a \omterm{def:g4:memory-and-storage:byte}{byte}, $1\,400 - 5 \times 256$.
The next book explains what happens when a number goes past 255.

\textbf{12.} \textbf{Result.} The multiplication of $a$ by $b$ takes $3b + 6$
steps: the same 3 more for each extra unit of $b$, so $3 \times 50$ takes 156
steps and $50 \times 3$ only 15. A \omterm{def:g4:memory-and-storage:byte}{one-byte} \omterm{def:g7:how-a-processor-runs-a-program:register}{register} cannot hold a product
above 255.
\end{solution}
